\chapter{Ownership, Lifetimes, and Garbage Collection at Language Boundaries}
\label{ch:ffi-ownership}

A native pointer can outlive the runtime object that originally justified it.
Keeping that object reachable fixes one problem, but a moving collector can still change its address.
Conversely, preventing movement does not necessarily keep an otherwise unreachable object alive.
An FFI must account for reachability, address stability, and release authority separately.

\section{State who may release each resource}
\label{sec:ch72-ownership}
\term{Foreign ownership} is the contract assigning responsibility for retaining and releasing a resource across the boundary.
A borrowed pointer grants temporary access without transferring release responsibility.
An owned result transfers a specified release obligation, often through an explicit destructor function.
Shared ownership adds a retention mechanism whose increments and decrements must agree across all participants.

Our native constructor copies bytes into a \code{BbBuffer} and publishes its pointer only after successful allocation and initialization.
The caller must eventually call \code{bb\_free}; it must not guess which allocator can free the result.
The allocation and destructor belong to the same native component, which also allows the representation to remain opaque.
Table~\ref{tab:ch72-resources} accounts for the resources used in the two adapters.

\begin{table}[htbp]
\caption{Every resource has both a retention mechanism and a release event.}
\label{tab:ch72-resources}
\centering
\begin{tabular}{lll}
\toprule
Resource & Retained by & Released at \\
\midrule
Python input export & buffer-protocol acquisition & end of constructor copy \\
Lua input string & argument on Lua stack & ordinary call completion \\
native byte buffer & capsule or full userdata owner & explicit close or finalization \\
callback function & synchronous call arguments & visitor return \\
prepared call metadata & address-stable stack plan & foreign call completion \\
\bottomrule
\end{tabular}
\end{table}

The C constructor leaves its output unchanged on failure.
The wrappers initialize their owner slots to null, so cleanup is valid even if construction fails before publication.
The native destructor accepts null, but it is not generally idempotent on a nonnull freed address.
Idempotent close belongs to the wrapper: it releases once and replaces its stored pointer with null.

\section{Separate roots, handles, and pins}
\label{sec:ch72-movement}
A \term{foreign root} keeps a managed object reachable for the collector.
A \term{stable handle} names an indirection slot whose target can be updated when an object moves.
A \term{pin} prevents relocation for a specified interval under the runtime's pinning rules.
These mechanisms solve related but distinct problems~\cite{jones-2023-garbage}.

Figure~\ref{fig:ch72-handle} shows a handle surviving relocation.
The foreign client keeps the handle, not a cached raw address from its previous lookup.
If it resolves the handle to a raw pointer and then allows a collection, it must reacquire the address afterward or enter a protocol guaranteeing stability during use.
The existence of the handle alone does not make every derived pointer permanent.

\begin{figure}[htbp]
\centering
\begin{tikzpicture}[node distance=13mm and 16mm]
\node[basic block] (handle) {rooted handle slot};
\node[basic block,below left=of handle] (old) {old object address};
\node[basic block,below right=of handle] (new) {new object address};
\draw[cfg edge,dashed] (handle)--node[left,font=\small]{before}(old);
\draw[cfg edge] (handle)--node[right,font=\small]{after}(new);
\draw[cfg edge] (old)--node[below,font=\small]{move}(new);
\end{tikzpicture}
\caption{A moving collector updates the rooted indirection while foreign code retains the handle.}
\label{fig:ch72-handle}
\end{figure}

The companion handle model stores a slot index and generation.
Releasing slot $(0,0)$ increments its generation; reusing the slot produces $(0,1)$.
A stale $(0,0)$ is rejected rather than resolving to the new owner's object.
Algorithm~\ref{alg:ch72-handle} gives the lookup rule.
Finite-width production generations must address wraparound, for example by retiring exhausted slots; an unbounded mathematical counter does not solve that implementation problem automatically.

\begin{algorithm}[htbp]
\caption{Resolve a sequential generation-tagged handle.}
\label{alg:ch72-handle}
\begin{pseudocode}
\Require slot table, handle $(i,g)$; table cannot change during this lookup/use
\Ensure current object or stale-handle rejection
\If{$i$ is outside the table} \State reject \EndIf
\State read slot generation $g'$ and object $o$
\If{$g\ne g'$ or the slot is empty} \State reject \EndIf
\State return $o$ under the caller's stability protocol
\end{pseudocode}
\end{algorithm}

Array lookup costs $O(1)$; the simple model finds a reusable slot by an $O(H)$ scan over $H$ slots.
Moving all live model objects costs $O(H)$ for their fixed-size records.
The model is sequential and does not establish lock-free lookup safety or concurrent reclamation.
Those require a protocol preventing release between validation and use, not merely an atomic generation comparison.

\section{Make close and finalization cooperate}
\label{sec:ch72-close}
A finalizer is a fallback for an unreachable wrapper, not a reliable deadline for releasing scarce resources.
Explicit close makes the event visible in program control flow.
After close, operations on a surviving wrapper must detect the closed state rather than using its former payload.
Figure~\ref{fig:ch72-close} shows the wrapper's two-state protocol.

\begin{figure}[htbp]
\centering
\begin{tikzpicture}[node distance=24mm]
\node[basic block] (live) {live payload};
\node[basic block,right=of live] (closed) {null payload};
\draw[cfg edge] (live)--node[above,font=\small]{close}(closed);
\draw[cfg edge] (closed) edge[loop below] node[below,font=\small]{close again}(closed);
\end{tikzpicture}
\caption{Wrapper-level close is idempotent because the pointer is cleared after release.}
\label{fig:ch72-close}
\end{figure}

Algorithm~\ref{alg:ch72-close} assumes serialized access and a destructor that neither throws nor reenters the wrapper.
If destruction can reenter, detach the payload from the owner before invoking the destructor and define how reentrant calls observe the closing state.
If operations can run concurrently, close must also stop new users and wait for existing users; clearing a pointer alone is insufficient.

\begin{algorithm}[htbp]
\caption{Close a serialized owner with a nonreentrant native destructor.}
\label{alg:ch72-close}
\begin{pseudocode}
\Require valid owner; no concurrent user; destructor accepts null and cannot reenter
\Ensure owner has no payload and repeated close is harmless
\State release the owner's payload through its matching native destructor
\State set the owner's payload pointer to null
\State return success
\end{pseudocode}
\end{algorithm}

The bookkeeping is $O(1)$ plus the destructor's cost.
In CPython, a capsule retains a small owner allocation until the capsule itself dies, even after its payload is closed.
In Lua, full userdata stores the owner directly; both \code{\_\_gc} and \code{\_\_close} use the same close routine.
The Lua adapter installs the metatable before acquiring the native payload, so any later allocation error still leaves a finalizable owner~\cite{ierusalimschy-2026-lua}.

\section{Avoid cycles that neither collector can see}
\label{sec:ch72-cycles}
Suppose a native registration retains a managed callback, and that callback captures the wrapper owning the registration.
The runtime sees the callback-to-wrapper edge, while the native table supplies the retention keeping the callback alive.
Without an explicit release event or integrated tracing, neither ordinary wrapper finalization nor dropping the user's last reference necessarily breaks the cycle.
An explicit registration token with close, a weak edge with defined failure behavior, or a tracing integration can resolve it.

Table~\ref{tab:ch72-policies} compares the principal lifetime choices.
Pinning may constrain collector movement and increase fragmentation; copying may consume bandwidth and duplicate storage.
The right policy follows the operation's duration, mutation needs, and runtime guarantees rather than a blanket preference for avoiding copies.

\begin{table}[htbp]
\caption{Lifetime mechanisms preserve different properties and impose different costs.}
\label{tab:ch72-policies}
\centering
\begin{tabular}{lll}
\toprule
Mechanism & Preserves & Main obligation \\
\midrule
root & reachability & eventually release the root \\
handle & identity across moves & resolve under a stability protocol \\
pin & stable address while pinned & retain liveness; eventually unpin \\
copy & independent storage & define snapshot or mutation rules \\
ownership transfer & designated release duty & use the specified destructor \\
\bottomrule
\end{tabular}
\end{table}

\section{Exercises}
\label{sec:ch72-exercises}
\begin{enumerate}
\item\label{ex:ch72-generation} Release handle $(0,0)$, reuse its slot for value $99$, and attempt both old and new lookups.
Give the results and explain what a generation-wrap policy must prevent.
\item\label{ex:ch72-pin} Explain why retaining an object does not necessarily stabilize its address and why a raw pin contract need not imply reachability.
\item\label{ex:ch72-close} Trace explicit close, repeated close, a read attempt, and finalization for one owner.
Count payload releases.
\item\label{ex:ch72-cycle} Draw the registration/callback/wrapper retention cycle in prose and identify one explicit operation that breaks it.
\end{enumerate}

\section{Further reading}
\label{sec:ch72-reading}
The second edition of the garbage-collection handbook explains foreign references, handles, and pinning in the context of moving collectors~\cite{jones-2023-garbage}.
The Python capsule and buffer APIs and Lua userdata/finalization APIs show concrete ownership mechanisms~\cite{python-2026-capi,ierusalimschy-2026-lua}.
